\documentclass[tikz,a4paper,11pt]{book}
%\documentclass[a4paper,twoside,11pt,titlepage]{book}
\usepackage{listings}
\usepackage[utf8]{inputenc}
\usepackage{csquotes}
\usepackage[spanish,es-noquoting]{babel}
\usepackage{forest}
\usepackage{xcolor}
\usepackage[export]{adjustbox}
\usepackage{caption}
\usepackage{float}
\usetikzlibrary{arrows.meta,shapes,positioning,trees}

% \usepackage[style=list, number=none]{glossary} %
%\usepackage{titlesec}
%\usepackage{pailatino}

\decimalpoint
\usepackage{dcolumn}
\newcolumntype{.}{D{.}{\esperiod}{-1}}
\makeatletter
\addto\shorthandsspanish{\let\esperiod\es@period@code}
\makeatother


%\usepackage[chapter]{algorithm}
\RequirePackage{verbatim}
%\RequirePackage[Glenn]{fncychap}
\usepackage{fancyhdr}
\usepackage{graphicx}
\usepackage{afterpage}
\usepackage{longtable}
\usepackage[pdfborder={000}]{hyperref} %referencia


% ********************************************************************
% Re-usable information
% ********************************************************************
\newcommand{\myTitle}{Chief\xspace}
\newcommand{\myDegree}{Grado en Ingeniería Informática\xspace}
\newcommand{\myName}{Sergio Cervilla Ortega\xspace}
\newcommand{\myProf}{Juan Julián Merelo Guervós\xspace}
\newcommand{\myFaculty}{Escuela Técnica Superior de Ingenierías Informática y de Telecomunicación\xspace}
\newcommand{\myFacultyShort}{E.T.S. de Ingenierías Informática y de Telecomunicación\xspace}
% \newcommand{\myDepartment}{Departamento de Arquitectura y Tecnología de Computadores\xspace}
\newcommand{\myUni}{\protect{Universidad de Granada}\xspace}
\newcommand{\myLocation}{Granada\xspace}
\newcommand{\myTime}{\today\xspace}
\newcommand{\myVersion}{Version 1.0\xspace}
\newcommand{\innerxsep}{5mm}

\hypersetup{
pdfauthor = {\myName (cervick13@gmail.com)},
pdftitle = {\myTitle},
pdfsubject = {},
pdfcreator = {LaTeX con el paquete ....},
pdfproducer = {pdflatex}
}

%\hyphenation{}


%\usepackage{doxygen/doxygen}
%\usepackage{pdfpages}
\usepackage{url}
\usepackage{hyperref}
\usepackage{colortbl,longtable}
\usepackage[stable]{footmisc}
%\usepackage{index}

%\makeindex
%\usepackage[style=long, cols=2,border=plain,toc=true,number=none]{glossary}
% \makeglossary

% Definición de comandos que me son tiles:
%\renewcommand{\indexname}{Índice alfabético}
%\renewcommand{\glossaryname}{Glosario}

\pagestyle{fancy}
\fancyhf{}
\fancyhead[LO]{\leftmark}
\fancyhead[RE]{\rightmark}
\fancyhead[RO,LE]{\textbf{\thepage}}


%Para conseguir que en las páginas en blanco no ponga cabecerass
\makeatletter
\def\clearpage{%
  \ifvmode
    \ifnum \@dbltopnum =\m@ne
      \ifdim \pagetotal <\topskip
        \hbox{}
      \fi
    \fi
  \fi
  \newpage
  \thispagestyle{empty}
  \write\m@ne{}
  \vbox{}
  \penalty -\@Mi
}
\makeatother


\usepackage{pdfpages}
\usepackage{fourier}
\begin{document}
\begin{titlepage}
\newlength{\centeroffset}
\setlength{\centeroffset}{-0.5\oddsidemargin}
\addtolength{\centeroffset}{0.5\evensidemargin}
\thispagestyle{empty}

\noindent\hspace*{\centeroffset}\begin{minipage}{\textwidth}

\centering
\includegraphics[width=0.9\textwidth]{imagenes/logo_ugr.jpg}\\[1.4cm]

\textsc{ \Large TRABAJO FIN DE GRADO\\[0.2cm]}
\textsc{ GRADO EN INGENIERIA INFORMATICA}\\[1cm]

{\Huge\bfseries Chief \\}
\noindent\rule[-1ex]{\textwidth}{3pt}\\[3.5ex]
{\large\bfseries Apoyo a la labor de la Policía Local }
\end{minipage}

\vspace{2.5cm}
\noindent\hspace*{\centeroffset}
\begin{minipage}{\textwidth}
\centering

\textbf{Autor}\\ {Sergio Cervilla Ortega}\\[2.5ex]
\textbf{Director}\\ {Juan Julián Merelo Guervós}\\[2cm]
\includegraphics[width=0.3\textwidth]{imagenes/etsiit_logo.png}\\[0.1cm]
\textsc{Escuela Técnica Superior de Ingenierías Informática y de Telecomunicación}\\
\textsc{---}\\
Granada, algún mes de 2018
\end{minipage}
\end{titlepage}

\chapter*{}
\thispagestyle{empty}

\begin{center}
{\large\bfseries Chief \\ Apoyo a la labor de la Policía Local}\\
\end{center}
\begin{center}
Sergio Cervilla Ortega\\
\end{center}

%\vspace{0.7cm}

\noindent{\textbf{Resumen}}\\

Lorem ipsum dolor sit amet, consectetur adipisicing elit, sed do eiusmod tempor incididunt ut labore et dolore magna aliqua. Ut enim ad minim veniam, quis nostrud exercitation ullamco laboris nisi ut aliquip ex ea commodo consequat. Duis aute irure dolor in reprehenderit in voluptate velit esse cillum dolore eu fugiat nulla pariatur. Excepteur sint occaecat cupidatat non proident, sunt in culpa qui officia deserunt mollit anim id est laborum.

\cleardoublepage

\thispagestyle{empty}

\noindent\rule[-1ex]{\textwidth}{2pt}\\[4.5ex]

D. \textbf{Juan Julián Merelo Guervós}, Profesor del Área de XXXX del Departamento de Arquitectura y Tecnología de Computadores de la Universidad de Granada.

\vspace{0.5cm}

\textbf{Informo:}

\vspace{0.5cm}

Que el presente trabajo, titulado \textit{\textbf{Chief, blah blah}},
ha sido realizado bajo mi supervisión por \textbf{Sergio Cervilla Ortega}, y autorizo la defensa de dicho trabajo ante el tribunal
que corresponda.

\vspace{0.5cm}

Y para que conste, expiden y firman el presente informe en Granada a X de mes de 2018.

\vspace{1cm}

\textbf{El director: }

\vspace{5cm}

\noindent \textbf{Juan Julián Merelo Guervós}

\frontmatter
\tableofcontents
\listoffigures
%\listoftables
%\mainmatter
%\setlength{\parskip}{5pt}

\chapter{Introducción}

Vivimos cada día en un mundo más conectado a Internet, es una época en la que se está dejando de 
lado el lápiz y el papel. Donde se está abriendo paso la tecnología. Estos nuevos medios, bien
usados, son capaces de facilitarnos en numerosas ocasiones las tareas del día a día. Convierte acciones
tediosas en simples gestos que, con el tiempo, acabamos automatizando y realizamos sin esfuerzo. 
Bajo dicha motivación, decidí crear un sistema informatizado para ayudar en las labores del cuerpo policial de Maracena.\\

Previamente, el cuerpo policial utilizaba un esquema bastante anticuado para desemplear sus funciones.
Por ejemplo, los registros de incidencias se iban apuntando manualmente y al final del turno
se reunían para poner en común todos los datos en un documento que posteriormente se guardaba.
Y al igual que en el ejemplo anterior, debían esperar a llegar a las dependencias para poder 
rellenar cualquier tipo de denuncia administrativa con los datos que habían apuntado en 
un papel para no olvidarlos.\\

Para ello, he diseñado una aplicación web que se encarga de gestionar una gran cantidad de 
modelos de denuncia y que permite que se rellenen de forma segura, a prueba de 
fallos y de una manera muy intuitiva para el usuario. Además de poder gestionar dichos documentos,
se ha creado un sistema de registro de incidencias completamente automatizado para que en 
ningún momento se pierda ningún dato y un gran número de funcionalidades más que se explicarán en 
posteriores puntos. El Ayuntamiento de Maracena será pionero en este ámbito, ya que han sido los 
primeros que han apostado por las nuevas tecnologías. Una gran apuesta, pero que traerá grandes
beneficios a la calidad del trabajo del cuerpo policial y a la seguridad de los ciudadanos. \\

Este Trabajo de Fin de Grado ha sido realizado completamente por mi, Sergio 
Cervilla Ortega durante el año 2018 y es libre bajo la licencia GNU General Public License v3.0 \cite{gplv3}\\ 

Ha sido desarrollado bajo el permiso, aceptación, revisión y ayuda del Ayuntamiento de Maracena, localidad en
la que resido a fecha de hoy.\\


\chapter{Tecnologías}

\subsection{Node.js}

Node.js es un entorno de ejecución para JavaScript construido con el motor de JavaScript
V8 de Chrome. En este proyecto se ha utilizado para elaborar el servidor de la aplicación.\\

Node.js tiene una característica que lo diferencia del resto de posibilidades para implementar 
el servidor y es que utiliza un modelo de operaciones E/S sin bloqueo y orientado a eventos 
asícronos. Esto nos ofrece la la posibilidad para que cuando se detecte una conexión se atienda
pero el resto del tiempo estará inactivo. Con este nuevo modelo conseguimos servidores altamente 
escalables de una manera muy sencilla, rápida y eficiente.\\

Este diseño contrasta directamente con el modelo con el modelo de concurrencia más común hoy en 
día en el que se usan hilos del Sistema Operativo. Las operaciones en red basadas en este modelo
son relativamente ineficientes además de complejas en su uso. Cabe mencionar que, aunque Node.js 
no utilice hilos del Sistema Operativo implique que no podamos aprovechar los múltiples cores
de nuestro sistema.\\

Otro punto a favor de Node.js es que no tendremos que preocuparnos por la posibilidad de un 
bloqueo en el proceso debido a que no existe dicha posibilidad por la naturaleza de dicha 
tecnología, la cual está basada en un bucle de eventos.

\subsection{MongoDB}

MongoDB es un sistema de base de datos multiplataforma y orientado a documentos de esquema libre.
Esto quiere decir que cada entrada puede tener un esquema de datos diferente al anterior, por
lo que podemos tener entradas que difieran en el número de atributos entre sí.\\

Debido a esta característica los datos no son guardados en registros, como en las bases SQL, 
sino que se guardan en documentos. Estos documentos son del tipo BSON, que es una representación
binaria de los archivos JSON.\\

Esta tecnología ha sido escrita en C++, por lo que está en contacto con el \textit{bare metal}. Esto
último es realmente importante ya que le permite optimizar los recursos y acceder a ellos de una 
manera extremadamente eficiente, con lo que consigue una velocidad muy alta en todas sus operaciones.\\

Cabe mencionar que MongoDB es una plataforma completamente distribuida, lo que nos proporciona 
un nuevo nivel de escalabilidad y disponibilidad. Esto quiere decir que a medida que nuestro desarrollo
crezca tanto en volumen de datos como en rendimiento MongoDB se escala automáticamente sin tiempos
de caída y sin cambios en nuestra aplicación.\\

Por último, cabe mencionar que la licencia de esta tecnología es GNU AGPL 3.0, por lo que se trata de un software de licencia
libre.

\subsection{Vagrant}

Vagrant es una herramienta para construir y administrar entornos de máquinas virtuales de una 
manera sencilla. Debido a que está diseñado para ser sencillo de utilizar y a que se centra en 
la automatización hace que utilicemos el menor tiempo posible en la preparación del entorno virtual
que necesitamos para comenzar a trabajar.\\

Nos proporciona un entorno muy rápido y simple de configurar, reproducible y portable utilizando
tecnología puntera en el sector de la virtualización como puede ser VirtualBox, VMWare, AWS, Docker... 


\subsection{Ansible}




\chapter{Elección de la plataforma de gestión documental}

Esta es una de las elecciones más significativas debido a que toda la gestión documental
y la manera en la que se realiza depende del software que utilicemos.
\\

En este caso en particular vamos a utilizar \textbf{Nuxeo} debido a que se adapta mejor a mis
necesidades que mi segunda opción, \textbf{Alfresco}.
\\

Los motivos por los que he elegido Nuxeo frente a Alfresco son los siguientes:

\begin{itemize}
	\item Alfresco sigue la filosofía \textit{Open Source} como método de trabajo, sino que la utilizan a modo de marketing
	 			ya que podemos comprobar que en su desarrollo solo intervienen trabajadores de dicha empresa a pesar de ser bastante
				más grande que Nuxeo.
	\item Alfresco tiene dos versiones, \textit{Community edition} y \textit{Commercial edition}. Esto no sería un problema
				si en la versión de la comunidad se incluyera \textbf{toda la funcionalidad}. Pero de esta manera, si quieres tener
				el software al completo deberás pagar la versión comercial. Sin embargo, Nuxeo nos proporciona toda la funcionalidad
				base tengas o no tengas contratada una suscripción.
	\item La tecnología que impulsa Nuxeo es también de código libre y actual.
	\item La \textit{API} de Nuxeo es mucho más completa.
	\item La documentación de Nuxeo es muy completa, fácil de acceder y de aplicar.
\end{itemize}

% \input{secciones/servicio_mensajeria}
% \input{secciones/planificacion}
% \input{secciones/implementacion}
% \input{secciones/conclusiones}

%
\nocite{*}
\bibliography{bibliografia}\addcontentsline{toc}{chapter}{Bibliografía}
\bibliographystyle{unsrt}


\end{document}
